\documentclass[11pt,a4paper]{article}
\usepackage[utf8]{inputenc}
\usepackage[T1]{fontenc}
\usepackage[catalan]{babel}
\usepackage{lmodern}
\usepackage[sfdefault,lining]{FiraSans}
\usepackage[a4paper,top=2.3cm,bottom=2.2cm,left=1.8cm,right=1.8cm,headheight=1.3cm,headsep=0.4cm,footskip=1cm]{geometry}
\usepackage{amsmath,amssymb}
\allowdisplaybreaks
\usepackage{icomma}
\usepackage{graphicx}
\usepackage[table]{xcolor}
\usepackage{tikz}
\usepackage{fontawesome5}
\usepackage{ccicons}
\usepackage{fancyhdr}
\usepackage{lastpage}
\usepackage{tasks}
\usepackage{needspace}
\usepackage[most]{tcolorbox}
\usepackage{array}
\usepackage[hidelinks,pdftitle={Solucions de la fitxa de la unitat 4: Problemes mètrics. 2n de Batxillerat MII},pdfsubject={Matemàtiques II},pdfkeywords={angles, distàncies, perpendicular comuna, projeccions, simètrics, àrees, volums, pla mediador, plans bisectors, esfera, PAU}]{hyperref}

\definecolor{navy}{HTML}{1B2748}
\definecolor{accent}{HTML}{E8603A}
\definecolor{teal}{HTML}{1F9E90}
\definecolor{soft}{HTML}{EEF1F8}
\definecolor{muted}{HTML}{5D6685}

\setlength{\parindent}{0pt}
\setlength{\parskip}{3pt}
\settasks{label=\alph*), label-format=\bfseries\color{navy}, label-width=1.3em, item-indent=1.8em, column-sep=0.6em, after-item-skip=5pt, before-skip=3pt, after-skip=4pt}

% ---------- capçalera i peu ----------
\newcommand{\llicencia}{\href{https://creativecommons.org/licenses/by-nc-sa/4.0/deed.ca}{\ccbyncsa}\ \ Llicència Creative Commons Reconeixement-NoComercial-CompartirIgual 4.0 Internacional (CC BY-NC-SA 4.0)}
\pagestyle{fancy}
\fancyhf{}
\fancyhead[L]{\raisebox{-0.35\height}{\includegraphics[height=0.95cm]{logo-gva.png}}}
\fancyhead[C]{\small\color{muted}Solucions de la fitxa · Unitat 4: Problemes mètrics · 2n Batxillerat MII}
\fancyhead[R]{\raisebox{-0.35\height}{\includegraphics[height=0.95cm]{logo-vermellar.png}}}
\fancyfoot[L]{\scriptsize\color{muted}\llicencia}
\fancyfoot[R]{\small\color{muted}\thepage}
\renewcommand{\headrulewidth}{0.4pt}
\renewcommand{\headrule}{\color{muted!40}\hrule height \headrulewidth}
\fancypagestyle{primera}{\fancyhf{}\renewcommand{\headrulewidth}{0pt}
  \fancyfoot[L]{\scriptsize\color{muted}\llicencia}
  \fancyfoot[R]{\small\color{muted}\thepage}}

% ---------- nivells i exercicis ----------
\newcommand{\estrelles}[1]{{\color{accent}\ifcase#1\or\faStar\faStar[regular]\faStar[regular]\or\faStar\faStar\faStar[regular]\or\faStar\faStar\faStar\fi}}
\newcounter{exer}
\newcommand{\ex}[1]{\par\needspace{4\baselineskip}\medskip\refstepcounter{exer}\noindent{\color{navy}\bfseries\theexer.}\ \ }

\newcommand{\seccio}[2]{\par\needspace{10\baselineskip}\bigskip
  \begin{tcolorbox}[colback=navy,colframe=navy,coltext=white,arc=2mm,boxrule=0pt,left=3mm,right=3mm,top=1.2mm,bottom=1.2mm]
  {\large\bfseries #1\quad #2}
  \end{tcolorbox}\vspace{-2pt}}


\setlength{\emergencystretch}{3em}  % les fórmules llargues (equacions de rectes, ternes) passen a la línia següent

% quadre de recordatori (només a la fitxa d'exercicis)
\newcommand{\recorda}[1]{\par\smallskip\begin{tcolorbox}[colback=accent!7,colframe=accent!55,arc=2mm,boxrule=0.5pt,left=2.5mm,right=2.5mm,top=1mm,bottom=1mm]\small{\color{accent}\faLightbulb[regular]}\ \textbf{Recorda:}\ #1\end{tcolorbox}}
% etiqueta dels exercicis de la PAU d'altres comunitats
\newcommand{\pau}[1]{{\color{teal}\small\bfseries\faUniversity\ PAU · #1.}\ }
\newcommand{\res}[1]{\boxed{#1}}\newcommand{\rest}[1]{\fbox{#1}}

% =====================================================================
% Capçalera FentMates de les fitxes: un color per a cada curs.
% Els colors són els de la pàgina inicial (LEVELS a assets/js/materials.js);
% el logo ve de _eines/logo/pdf/ (python3 _eines/logo/genera.py).
%
% Ús, al final del preàmbul de la unitat:
%   % =====================================================================
% Capçalera FentMates de les fitxes: un color per a cada curs.
% Els colors són els de la pàgina inicial (LEVELS a assets/js/materials.js);
% el logo ve de _eines/logo/pdf/ (python3 _eines/logo/genera.py).
%
% Ús, al final del preàmbul de la unitat:
%   % =====================================================================
% Capçalera FentMates de les fitxes: un color per a cada curs.
% Els colors són els de la pàgina inicial (LEVELS a assets/js/materials.js);
% el logo ve de _eines/logo/pdf/ (python3 _eines/logo/genera.py).
%
% Ús, al final del preàmbul de la unitat:
%   \input{../comu/fentmates.tex}
%   \fmcurs{1eso}
% i, al començament del document:
%   \capcalera{Matemàtiques}{1r ESO}{Fitxa d'exercicis · Unitat 2}{Divisibilitat}
% Opcional abans de \input: \def\fmneed{14} (espai mínim abans de cada secció).
% =====================================================================
\makeatletter
\def\fm@logodir{../../logo/pdf}
\providecommand{\fmneed}{10}
\def\fm@cursnom{}

% clau del curs, color i color del text damunt del color
\newcommand{\fm@curs}[3]{\@namedef{fm@color@#1}{#2}\@namedef{fm@text@#1}{#3}}
\fm@curs{1eso}{E6B84A}{navy}
\fm@curs{2eso}{3B82F6}{white}
\fm@curs{3eso}{FF7A50}{white}
\fm@curs{4eso-a}{A78BFA}{white}
\fm@curs{4eso-b}{6D28D9}{white}
\fm@curs{1batx-mi}{9B1C1C}{white}
\fm@curs{1batx-ccss}{EF6B6B}{white}
\fm@curs{2batx-mii}{1E6B47}{white}
\fm@curs{2batx-ccss}{2BB3A3}{white}

\newcommand{\fmcurs}[1]{%
  \@ifundefined{fm@color@#1}{\PackageError{fentmates}{Curs desconegut: #1}{}}{}%
  \edef\fm@tmp{\noexpand\definecolor{curs}{HTML}{\@nameuse{fm@color@#1}}%
               \noexpand\colorlet{curstext}{\@nameuse{fm@text@#1}}}\fm@tmp
  \colorlet{cursfosc}{curs!72!black}%
  \def\fm@logo{\fm@logodir/vertical-#1.pdf}%
  \def\fm@marca{\fm@logodir/marca-#1.pdf}}

% ---------- capçalera de la primera pàgina ----------
% #1 matèria (petit), #2 curs, #3 tipus de document i unitat, #4 títol
\newcommand{\capcalera}[4]{%
  \gdef\fm@cursnom{#2}%
  \par\noindent
  \begin{tcolorbox}[colback=curs,colframe=curs,boxrule=0pt,arc=2.2mm,width=3.5cm,height=2.15cm,
      valign=center,left=3mm,right=2mm,top=1mm,bottom=1mm,nobeforeafter,box align=center]
    \color{curstext}{\footnotesize #1}\par\vspace{1pt}{\fontsize{19}{21}\selectfont\bfseries #2\par}%
  \end{tcolorbox}\hspace{5mm}%
  \begin{minipage}[c]{\dimexpr\linewidth-3.5cm-5mm-2.9cm\relax}
    {\color{cursfosc}\small\bfseries\MakeUppercase{#3}}\par\vspace{2pt}
    {\color{navy}\fontsize{20}{23}\selectfont\bfseries #4\par}
  \end{minipage}\hfill
  \begin{minipage}[c]{2.6cm}\raggedleft\includegraphics[width=2.6cm]{\fm@logo}\end{minipage}\par
  \vspace{2pt}}

% ---------- barres de secció del color del curs ----------
\renewcommand{\seccio}[2]{\par\needspace{\fmneed\baselineskip}\bigskip
  \begin{tcolorbox}[colback=curs,colframe=curs,coltext=curstext,arc=2mm,boxrule=0pt,left=3mm,right=3mm,top=1.2mm,bottom=1.2mm]
  {\large\bfseries #1\quad #2}
  \end{tcolorbox}\vspace{-2pt}}

% ---------- estrelles de dificultat del color del curs (a Batxillerat no se'n posen) ----------
\@ifundefined{estrelles}{}{%
  \renewcommand{\estrelles}[1]{{\color{curs}\ifcase#1\or\faStar\faStar[regular]\faStar[regular]\or\faStar\faStar\faStar[regular]\or\faStar\faStar\faStar\fi}}}

% ---------- capçalera i peu de les pàgines ----------
% primera pàgina: logos institucionals al peu
\fancypagestyle{primera}{\fancyhf{}\renewcommand{\headrulewidth}{0pt}
  \fancyfoot[L]{\raisebox{-0.32\height}{\includegraphics[height=0.78cm]{logo-gva.png}}\hspace{5mm}%
    \raisebox{-0.32\height}{\includegraphics[height=0.78cm]{logo-vermellar.png}}}
  \fancyfoot[R]{\scriptsize\color{muted}\href{https://creativecommons.org/licenses/by-nc-sa/4.0/deed.ca}{\ccbyncsa}\ \ CC BY-NC-SA 4.0 · fentmates.org\qquad\small\thepage}}
% la resta: el símbol a l'esquerra i el curs a la dreta, amb el títol al centre com abans
\fancyhead[L]{\raisebox{-0.3\height}{\includegraphics[height=0.62cm]{\fm@marca}}}
\fancyhead[R]{\small\bfseries\color{cursfosc}\fm@cursnom}
\renewcommand{\headrule}{\color{curs!45}\hrule height \headrulewidth}
\makeatother

%   \fmcurs{1eso}
% i, al començament del document:
%   \capcalera{Matemàtiques}{1r ESO}{Fitxa d'exercicis · Unitat 2}{Divisibilitat}
% Opcional abans de \input: \def\fmneed{14} (espai mínim abans de cada secció).
% =====================================================================
\makeatletter
\def\fm@logodir{../../logo/pdf}
\providecommand{\fmneed}{10}
\def\fm@cursnom{}

% clau del curs, color i color del text damunt del color
\newcommand{\fm@curs}[3]{\@namedef{fm@color@#1}{#2}\@namedef{fm@text@#1}{#3}}
\fm@curs{1eso}{E6B84A}{navy}
\fm@curs{2eso}{3B82F6}{white}
\fm@curs{3eso}{FF7A50}{white}
\fm@curs{4eso-a}{A78BFA}{white}
\fm@curs{4eso-b}{6D28D9}{white}
\fm@curs{1batx-mi}{9B1C1C}{white}
\fm@curs{1batx-ccss}{EF6B6B}{white}
\fm@curs{2batx-mii}{1E6B47}{white}
\fm@curs{2batx-ccss}{2BB3A3}{white}

\newcommand{\fmcurs}[1]{%
  \@ifundefined{fm@color@#1}{\PackageError{fentmates}{Curs desconegut: #1}{}}{}%
  \edef\fm@tmp{\noexpand\definecolor{curs}{HTML}{\@nameuse{fm@color@#1}}%
               \noexpand\colorlet{curstext}{\@nameuse{fm@text@#1}}}\fm@tmp
  \colorlet{cursfosc}{curs!72!black}%
  \def\fm@logo{\fm@logodir/vertical-#1.pdf}%
  \def\fm@marca{\fm@logodir/marca-#1.pdf}}

% ---------- capçalera de la primera pàgina ----------
% #1 matèria (petit), #2 curs, #3 tipus de document i unitat, #4 títol
\newcommand{\capcalera}[4]{%
  \gdef\fm@cursnom{#2}%
  \par\noindent
  \begin{tcolorbox}[colback=curs,colframe=curs,boxrule=0pt,arc=2.2mm,width=3.5cm,height=2.15cm,
      valign=center,left=3mm,right=2mm,top=1mm,bottom=1mm,nobeforeafter,box align=center]
    \color{curstext}{\footnotesize #1}\par\vspace{1pt}{\fontsize{19}{21}\selectfont\bfseries #2\par}%
  \end{tcolorbox}\hspace{5mm}%
  \begin{minipage}[c]{\dimexpr\linewidth-3.5cm-5mm-2.9cm\relax}
    {\color{cursfosc}\small\bfseries\MakeUppercase{#3}}\par\vspace{2pt}
    {\color{navy}\fontsize{20}{23}\selectfont\bfseries #4\par}
  \end{minipage}\hfill
  \begin{minipage}[c]{2.6cm}\raggedleft\includegraphics[width=2.6cm]{\fm@logo}\end{minipage}\par
  \vspace{2pt}}

% ---------- barres de secció del color del curs ----------
\renewcommand{\seccio}[2]{\par\needspace{\fmneed\baselineskip}\bigskip
  \begin{tcolorbox}[colback=curs,colframe=curs,coltext=curstext,arc=2mm,boxrule=0pt,left=3mm,right=3mm,top=1.2mm,bottom=1.2mm]
  {\large\bfseries #1\quad #2}
  \end{tcolorbox}\vspace{-2pt}}

% ---------- estrelles de dificultat del color del curs (a Batxillerat no se'n posen) ----------
\@ifundefined{estrelles}{}{%
  \renewcommand{\estrelles}[1]{{\color{curs}\ifcase#1\or\faStar\faStar[regular]\faStar[regular]\or\faStar\faStar\faStar[regular]\or\faStar\faStar\faStar\fi}}}

% ---------- capçalera i peu de les pàgines ----------
% primera pàgina: logos institucionals al peu
\fancypagestyle{primera}{\fancyhf{}\renewcommand{\headrulewidth}{0pt}
  \fancyfoot[L]{\raisebox{-0.32\height}{\includegraphics[height=0.78cm]{logo-gva.png}}\hspace{5mm}%
    \raisebox{-0.32\height}{\includegraphics[height=0.78cm]{logo-vermellar.png}}}
  \fancyfoot[R]{\scriptsize\color{muted}\href{https://creativecommons.org/licenses/by-nc-sa/4.0/deed.ca}{\ccbyncsa}\ \ CC BY-NC-SA 4.0 · fentmates.org\qquad\small\thepage}}
% la resta: el símbol a l'esquerra i el curs a la dreta, amb el títol al centre com abans
\fancyhead[L]{\raisebox{-0.3\height}{\includegraphics[height=0.62cm]{\fm@marca}}}
\fancyhead[R]{\small\bfseries\color{cursfosc}\fm@cursnom}
\renewcommand{\headrule}{\color{curs!45}\hrule height \headrulewidth}
\makeatother

%   \fmcurs{1eso}
% i, al començament del document:
%   \capcalera{Matemàtiques}{1r ESO}{Fitxa d'exercicis · Unitat 2}{Divisibilitat}
% Opcional abans de \input: \def\fmneed{14} (espai mínim abans de cada secció).
% =====================================================================
\makeatletter
\def\fm@logodir{../../logo/pdf}
\providecommand{\fmneed}{10}
\def\fm@cursnom{}

% clau del curs, color i color del text damunt del color
\newcommand{\fm@curs}[3]{\@namedef{fm@color@#1}{#2}\@namedef{fm@text@#1}{#3}}
\fm@curs{1eso}{E6B84A}{navy}
\fm@curs{2eso}{3B82F6}{white}
\fm@curs{3eso}{FF7A50}{white}
\fm@curs{4eso-a}{A78BFA}{white}
\fm@curs{4eso-b}{6D28D9}{white}
\fm@curs{1batx-mi}{9B1C1C}{white}
\fm@curs{1batx-ccss}{EF6B6B}{white}
\fm@curs{2batx-mii}{1E6B47}{white}
\fm@curs{2batx-ccss}{2BB3A3}{white}

\newcommand{\fmcurs}[1]{%
  \@ifundefined{fm@color@#1}{\PackageError{fentmates}{Curs desconegut: #1}{}}{}%
  \edef\fm@tmp{\noexpand\definecolor{curs}{HTML}{\@nameuse{fm@color@#1}}%
               \noexpand\colorlet{curstext}{\@nameuse{fm@text@#1}}}\fm@tmp
  \colorlet{cursfosc}{curs!72!black}%
  \def\fm@logo{\fm@logodir/vertical-#1.pdf}%
  \def\fm@marca{\fm@logodir/marca-#1.pdf}}

% ---------- capçalera de la primera pàgina ----------
% #1 matèria (petit), #2 curs, #3 tipus de document i unitat, #4 títol
\newcommand{\capcalera}[4]{%
  \gdef\fm@cursnom{#2}%
  \par\noindent
  \begin{tcolorbox}[colback=curs,colframe=curs,boxrule=0pt,arc=2.2mm,width=3.5cm,height=2.15cm,
      valign=center,left=3mm,right=2mm,top=1mm,bottom=1mm,nobeforeafter,box align=center]
    \color{curstext}{\footnotesize #1}\par\vspace{1pt}{\fontsize{19}{21}\selectfont\bfseries #2\par}%
  \end{tcolorbox}\hspace{5mm}%
  \begin{minipage}[c]{\dimexpr\linewidth-3.5cm-5mm-2.9cm\relax}
    {\color{cursfosc}\small\bfseries\MakeUppercase{#3}}\par\vspace{2pt}
    {\color{navy}\fontsize{20}{23}\selectfont\bfseries #4\par}
  \end{minipage}\hfill
  \begin{minipage}[c]{2.6cm}\raggedleft\includegraphics[width=2.6cm]{\fm@logo}\end{minipage}\par
  \vspace{2pt}}

% ---------- barres de secció del color del curs ----------
\renewcommand{\seccio}[2]{\par\needspace{\fmneed\baselineskip}\bigskip
  \begin{tcolorbox}[colback=curs,colframe=curs,coltext=curstext,arc=2mm,boxrule=0pt,left=3mm,right=3mm,top=1.2mm,bottom=1.2mm]
  {\large\bfseries #1\quad #2}
  \end{tcolorbox}\vspace{-2pt}}

% ---------- estrelles de dificultat del color del curs (a Batxillerat no se'n posen) ----------
\@ifundefined{estrelles}{}{%
  \renewcommand{\estrelles}[1]{{\color{curs}\ifcase#1\or\faStar\faStar[regular]\faStar[regular]\or\faStar\faStar\faStar[regular]\or\faStar\faStar\faStar\fi}}}

% ---------- capçalera i peu de les pàgines ----------
% primera pàgina: logos institucionals al peu
\fancypagestyle{primera}{\fancyhf{}\renewcommand{\headrulewidth}{0pt}
  \fancyfoot[L]{\raisebox{-0.32\height}{\includegraphics[height=0.78cm]{logo-gva.png}}\hspace{5mm}%
    \raisebox{-0.32\height}{\includegraphics[height=0.78cm]{logo-vermellar.png}}}
  \fancyfoot[R]{\scriptsize\color{muted}\href{https://creativecommons.org/licenses/by-nc-sa/4.0/deed.ca}{\ccbyncsa}\ \ CC BY-NC-SA 4.0 · fentmates.org\qquad\small\thepage}}
% la resta: el símbol a l'esquerra i el curs a la dreta, amb el títol al centre com abans
\fancyhead[L]{\raisebox{-0.3\height}{\includegraphics[height=0.62cm]{\fm@marca}}}
\fancyhead[R]{\small\bfseries\color{cursfosc}\fm@cursnom}
\renewcommand{\headrule}{\color{curs!45}\hrule height \headrulewidth}
\makeatother

\fmcurs{2batx-mii}

\begin{document}
\thispagestyle{primera}

\capcalera{Matemàtiques II}{2n Batx}{Solucions de la fitxa · Unitat 4}{\raggedright Problemes mètrics}
\vspace{8pt}

\begin{tcolorbox}[colback=soft,colframe=soft,arc=2mm,left=3mm,right=3mm,top=1.5mm,bottom=1.5mm]
\small Intenta resoldre cada exercici abans de mirar-ne la solució. Les distàncies i els angles es donen exactes (amb arrels simplificades) i, quan cal, aproximats amb dues xifres decimals. Si una recta o un pla de la teua solució s'escriu d'una altra manera, comprova que és el mateix (un punt comú i vectors proporcionals). Pots comprovar els resultats amb el laboratori mètric dels apunts interactius.
\end{tcolorbox}

% =====================================================
\seccio{1}{Angles entre rectes i plans}

\ex{1} \begin{tasks}(1)
  \task \mbox{$\cos\alpha=\dfrac{|8+0-6|}{3\cdot5}=\dfrac{2}{15}$} \ $\Rightarrow$ \ $\res{\alpha\approx82{,}34^\circ}$
  \task $\vec d_r=(1,\,1,\,0)\times(0,\,1,\,-1)=(-1,\,1,\,1)$; \ \mbox{$\cos\alpha=\dfrac{|-1+1+1|}{\sqrt3\cdot\sqrt3}=\dfrac13$} \ $\Rightarrow$ \ $\res{\alpha\approx70{,}53^\circ}$
\end{tasks}

\ex{1} \begin{tasks}(1)
  \task \mbox{$\sin\alpha=\dfrac{|1-1+2|}{\sqrt6\cdot\sqrt3}=\dfrac{2}{\sqrt{18}}=\dfrac{\sqrt2}{3}$} \ $\Rightarrow$ \ $\res{\alpha\approx28{,}13^\circ}$
  \task $\vec n_2=(2,\,-2,\,4)=2\,\vec d$: la recta és perpendicular al pla, \ $\res{\alpha=90^\circ}$
\end{tasks}

\ex{1} \begin{tasks}(1)
  \task \mbox{$\cos\alpha=\dfrac{1}{\sqrt2\sqrt2}=\dfrac12$}: \ $\res{\alpha=60^\circ}$
  \task $\vec n_{OXY}=(0,\,0,\,1)$: \ \mbox{$\cos\alpha=\dfrac23$}, \ $\res{\alpha\approx48{,}19^\circ}$
\end{tasks}

\ex{2} Amb $\vec i$, $\vec j$, $\vec k$: \ \mbox{$\cos\alpha=\dfrac23$}, \ \mbox{$\cos\beta=\dfrac13$}, \ \mbox{$\cos\gamma=\dfrac23$} \ $\Rightarrow$ \ $\res{\alpha=\gamma\approx48{,}19^\circ,\ \beta\approx70{,}53^\circ}$ \ (són els cosinus directors: $\cos^2\alpha+\cos^2\beta+\cos^2\gamma=1$).

\ex{2} \mbox{$\sin30^\circ=\dfrac{|1|}{\sqrt{2+m^2}}=\dfrac12$} \ $\Rightarrow$ \ $2+m^2=4$ \ $\Rightarrow$ \ $\res{m=\pm\sqrt2}$

\ex{2} $\vec n_1\cdot\vec n_2=2-k+k=2\neq0$: \ \rest{mai no són perpendiculars}. Per a ser paral·lels caldria \mbox{$\dfrac12=\dfrac{k}{-1}=\dfrac1k$}, és a dir, $k=-\tfrac12$ i $k=2$ alhora: \ \rest{tampoc no poden ser paral·lels}.

\ex{2} \pau{Canàries, extraordinària 2026}\begin{tasks}(1)
  \task $r:(-21+11t,\ 3-t,\ 5-2t)$; \ en $\pi$: $4t-2=10$, $t=3$: \ $\res{Q(12,\,0,\,-1)}$
  \task $\vec d=(11,\,-1,\,-2)$, $\vec n=(1,\,3,\,2)$: \ \mbox{$\sin\alpha=\dfrac{|11-3-4|}{\sqrt{126}\,\sqrt{14}}=\dfrac{4}{42}=\dfrac{2}{21}$} \ $\Rightarrow$ \ $\res{\alpha\approx5{,}47^\circ}$
\end{tasks}

\ex{2} \pau{Castella-la Manxa, ordinària 2024}$\vec u\cdot\vec v=2a-1=|\vec u|\,|\vec v|\cos60^\circ=\tfrac12\sqrt{1+2a^2}\,\sqrt5$. \ Cal $2a-1>0$; elevant al quadrat: $6a^2-16a-1=0$, \ \mbox{$a=\dfrac{8\pm\sqrt{70}}{6}$}. \ Només serveix \ \mbox{$\res{a=\dfrac{8+\sqrt{70}}{6}\approx2{,}73}$} \ (l'altra dona $120^\circ$).

\ex{2} \pau{País Basc, extraordinària 2025}\begin{tasks}(1)
  \task $\vec n\parallel\overrightarrow{AB}\times\overrightarrow{AC}=(-16,\,-16,\,8)$, \ $\vec n=(2,\,2,\,-1)$, \ $\vec d=(1,\,6,\,8)$: \ \mbox{$\sin\alpha=\dfrac{6}{3\sqrt{101}}=\dfrac{2}{\sqrt{101}}$} \ $\Rightarrow$ \ $\res{\alpha\approx11^\circ\,28'\,44''}$
  \task $\pi:2x+2y-z-6=0$; \ es tallen ($\vec d\cdot\vec n\neq0$) i $D$ ja és del pla: \ $\res{D(6,\,-5,\,-4)}$
\end{tasks}

\ex{2} \pau{Canàries, extraordinària 2024}\begin{tasks}(1)
  \task $|\vec u\cdot\vec v|\le|\vec u|\,|\vec v|=6<8$ (seria $\cos\alpha=\tfrac43>1$): \ \rest{no poden existir}.
  \task $\vec d\parallel(2,\,-3,\,-1)\times(1,\,1,\,2)=(-5,\,-5,\,5)$, $\vec d=(1,\,1,\,-1)$. \ Punt: $\res{P(1,\,0,\,-2)}$. \ \mbox{$\sin\alpha=\dfrac{2}{\sqrt3\sqrt{14}}$} \ $\Rightarrow$ \ $\res{\alpha\approx17{,}98^\circ}$
\end{tasks}

\ex{2} \pau{La Rioja, ordinària 2024}$\vec d_r=(1,\,1,\,1)\times(4,\,-1,\,1)=(2,\,3,\,-5)$. \ $r\cap\pi=\left(\tfrac25,\,-\tfrac75,\,0\right)$. \ $\vec d_s=\vec n\times\vec d_r=(1,\,1,\,1)$: \ $\res{s:x-\tfrac25=y+\tfrac75=z}$


% =====================================================
\seccio{2}{Distància entre punts i d'un punt a un pla}

\ex{1} \begin{tasks}(1)
  \task \mbox{$\dfrac{|6+6+6+1|}{\sqrt{9+36+4}}=\res{\dfrac{19}{7}}$}
  \task El pla és $3x+2y+z-6=0$: \ \mbox{$\dfrac{6}{\sqrt{14}}=\res{\dfrac{3\sqrt{14}}{7}}$}
  \task $|(2,\,-3,\,6)|=\sqrt{49}=\res{7}$
\end{tasks}

\ex{1} Recta per $P$ perpendicular: $(3+t,\ 3+t,\ -1-t)$; \ en el pla, $7+3t=1$, $t=-2$: \ $\res{Q(1,\,1,\,1)}$. \ $|\overrightarrow{PQ}|=|(-2,\,-2,\,2)|=2\sqrt3$ \ i \ \mbox{$d(P,\pi)=\dfrac{6}{\sqrt3}=2\sqrt3$} \ \checkmark

\ex{2} $(0,\,0,\,t)$: \ \mbox{$\dfrac{|-2t-1|}{3}=3$} \ $\Rightarrow$ \ $-2t-1=\pm9$: \ $\res{(0,\,0,\,-5)\ \text{i}\ (0,\,0,\,4)}$

\ex{2} \mbox{$\dfrac{|1+2a+2-3|}{3}=\dfrac{|2a|}{3}=2$} \ $\Rightarrow$ \ $\res{a=\pm3}$

\ex{2} \pau{Astúries, ordinària 2026}\begin{tasks}(1)
  \task Cal un director perpendicular a $\vec n=(1,\,2,\,3)$; per exemple, $\res{(x,\,y,\,z)=(0,\,-1,\,-2)+\lambda\,(1,\,1,\,-1)}$. \rest{No és única} (n'hi ha infinites).
  \task $(t,\ -1+2t,\ -2+3t)$: $14t-14=0$, $t=1$: \ $\res{Q(1,\,1,\,1)}$, \ $d(P,Q)=\res{\sqrt{14}}$
\end{tasks}

\ex{2} \pau{Catalunya, extraordinària 2024}\begin{tasks}(1)
  \task $(1+t,\ 3+2t,\ -2t)$: $14+9t=0$: \ $\res{\left(-\tfrac59,\,-\tfrac19,\,\tfrac{28}{9}\right)}$
  \task \mbox{$d=\dfrac{|1+6+7|}{3}=\res{\dfrac{14}{3}}$}
  \task $x+2y-2z+k=0$ amb $|7+k|=14$, $k\neq7$: \ $\res{x+2y-2z-21=0}$
\end{tasks}

\ex{2} \pau{Galícia, ordinària 2026}\begin{tasks}(1)
  \task $\overrightarrow{PQ}\times\overrightarrow{PR}=(-1,\,1,\,-1)\times(0,\,1,\,0)=(1,\,0,\,-1)\neq\vec0$: \ $\res{\pi:x-z=0}$
  \task \mbox{$\dfrac{|a|}{\sqrt2}=\sqrt2$}: \ $\res{a=\pm2}$
\end{tasks}

\ex{2} \pau{Astúries, extraordinària 2026}\begin{tasks}(1)
  \task $2x+2y+z+k=0$ amb \mbox{$\dfrac{|k|}{3}=2$}: \ \rest{n'hi ha dos}, \ $\res{2x+2y+z\pm6=0}$
  \task Projecció de $O$: \ $\res{P\left(\tfrac{10}{9},\,\tfrac{10}{9},\,\tfrac59\right)}$
  \task \mbox{$|\overrightarrow{OP}\times\overrightarrow{OQ}|=\left|\left(\tfrac53,\,-\tfrac53,\,0\right)\right|=\res{\dfrac{5\sqrt2}{3}}$}
\end{tasks}

\ex{2} \pau{Andalusia, ordinària 2025}\begin{tasks}(1)
  \task $[\overrightarrow{OA},\overrightarrow{OB},\overrightarrow{OC}]=4m-2$; \ \mbox{$\dfrac{|4m-2|}{6}=3$}: \ $\res{m=5\ \text{o}\ m=-4}$
  \task Pla $ABC$: $2x-z-2=0$; \ \mbox{$d(O,\pi)=\dfrac{2}{\sqrt5}=\res{\dfrac{2\sqrt5}{5}}$}
\end{tasks}


% =====================================================
\seccio{3}{Distància d'un punt a una recta}

\ex{1} \begin{tasks}(1)
  \task $A=O$, $\vec d=(1,\,1,\,1)$: \ $\overrightarrow{OP}\times\vec d=(-1,\,2,\,-1)$, \ \mbox{$d=\dfrac{\sqrt6}{\sqrt3}=\res{\sqrt2}$}
  \task $\vec d=(1,\,1,\,0)$, $A(0,\,0,\,2)$: \ $\overrightarrow{AP}\times\vec d=(3,\,-3,\,-1)$, \ \mbox{$d=\dfrac{\sqrt{19}}{\sqrt2}=\res{\dfrac{\sqrt{38}}{2}}$}
\end{tasks}

\ex{2} $\overrightarrow{QP}\cdot\vec d=0$: \ $(2-\lambda)+(1-\lambda)+(3-\lambda)=0$, $\lambda=2$: \ $\res{Q(3,\,2,\,3)}$. \ $d(P,r)=|\overrightarrow{QP}|=|(0,\,-1,\,1)|=\res{\sqrt2}$. \ Recta: \ $\res{(x,\,y,\,z)=(3,\,1,\,4)+\mu\,(0,\,1,\,-1)}$

\ex{2} $d=|\lambda|\cdot|(2,\,1,\,2)|=3|\lambda|=3$ \ $\Rightarrow$ \ $\lambda=\pm1$: \ $\res{(3,\,2,\,2)\ \text{i}\ (-1,\,0,\,-2)}$

\ex{2} És $d(C,r_{AB})$: \ $\overrightarrow{AB}=(-1,\,1,\,0)$, \ $\overrightarrow{AC}\times\overrightarrow{AB}=(-1,\,-1,\,1)$ \ $\Rightarrow$ \ \mbox{$h=\dfrac{\sqrt3}{\sqrt2}=\res{\dfrac{\sqrt6}{2}}$}

\ex{2} \pau{La Rioja, extraordinària 2026}\begin{tasks}(1)
  \task $\res{3x+y-z=0}$
  \task $A(0,\,0,\,11)$, $\overrightarrow{AP}\times\vec d=(0,\,1,\,-10)\times(3,\,1,\,-1)=(9,\,-30,\,-3)$: \ \mbox{$d=\dfrac{\sqrt{990}}{\sqrt{11}}=\sqrt{90}=\res{3\sqrt{10}}$}
\end{tasks}

\ex{2} \pau{Cantàbria, ordinària 2025}\begin{tasks}(1)
  \task $r:(2+t,\ 1+2t,\ t)$; \ $6t^2+2=14$, $t=\pm\sqrt2$: \ $\res{\left(2\pm\sqrt2,\ 1\pm2\sqrt2,\ \pm\sqrt2\right)}$
  \task $\res{x-y+z-1=0}$
  \task $\overrightarrow{AP}\times\vec d=(-1,\,0,\,1)\times(1,\,2,\,1)=(-2,\,2,\,-2)$: \ \mbox{$d=\dfrac{\sqrt{12}}{\sqrt6}=\res{\sqrt2}$}
\end{tasks}

\ex{2} \pau{Extremadura, extraordinària 2025}\begin{tasks}(1)
  \task $\res{4x-3y-7=0}$
  \task $\overrightarrow{AP}\times\vec d=(1,\,-2,\,1)\times(-4,\,3,\,0)=(-3,\,-4,\,-5)$: \ \mbox{$d=\dfrac{\sqrt{50}}{5}=\res{\sqrt2}$}
  \task \mbox{$\sin\alpha=\dfrac{|-4-6|}{5\sqrt{14}}=\dfrac{2}{\sqrt{14}}$} \ $\Rightarrow$ \ $\res{\alpha\approx32{,}31^\circ}$
\end{tasks}

\ex{2} \pau{Cantàbria, ordinària 2026}\begin{tasks}(1)
  \task $\res{4x+6y-3z+12=0}$
  \task $\overrightarrow{AP}\times\overrightarrow{AB}=(1,\,1,\,-3)\times(-3,\,0,\,-4)=(-4,\,13,\,3)$: \ \mbox{$d=\dfrac{\sqrt{194}}{5}\approx\res{2{,}79}$}
  \task $V=\tfrac16|{-24}|=\res{4}$
\end{tasks}


% =====================================================
\seccio{4}{Distància entre dues rectes}

\ex{1} Directors oposats i $(2,\,0,\,3)\notin r$ ($\overrightarrow{PQ}=(1,\,-1,\,3)$ no és proporcional a $(1,\,-1,\,1)$). \ $\overrightarrow{PQ}\times\vec d_r=(2,\,2,\,0)$ \ $\Rightarrow$ \ \mbox{$d=\dfrac{\sqrt8}{\sqrt3}=\res{\dfrac{2\sqrt6}{3}}$}

\ex{2} \label{ex:pc}$r$: $O$, $(1,\,1,\,0)$; \ $s$: $(1,\,0,\,0)$, $(0,\,1,\,1)$. \ $\vec d_r\times\vec d_s=(1,\,-1,\,1)$, \ $[\overrightarrow{PQ},\vec d_r,\vec d_s]=1\neq0$: s'encreuen; \ \mbox{$d=\dfrac{1}{\sqrt3}=\res{\dfrac{\sqrt3}{3}}$}

\ex{2} $\vec d_r\times\vec d_s=(1,\,0,\,0)\times(1,\,1,\,0)=(0,\,0,\,1)$, \ $[\overrightarrow{OQ},\vec d_r,\vec d_s]=(0,\,1,\,a)\cdot(0,\,0,\,1)=a$: \ $d=|a|=2$ \ $\Rightarrow$ \ $\res{a=\pm2}$

\ex{2} \pau{País Basc, ordinària 2025}\begin{tasks}(1)
  \task $\vec d_r=(5,\,-4,\,-2)$, $\vec d_s=(4,\,-4,\,0)$, $[\overrightarrow{AC},\vec d_r,\vec d_s]=-32\neq0$: \ \rest{s'encreuen}
  \task $|\vec d_r\times\vec d_s|=|(-8,\,-8,\,-4)|=12$: \ \mbox{$d=\dfrac{32}{12}=\res{\dfrac83}$}
\end{tasks}

\ex{2} \pau{La Rioja, extraordinària 2025}$\vec d_r\times\vec d_s=(2,\,2,\,-1)$, $\overrightarrow{PQ}=(-3,\,3,\,-9)$, $[\overrightarrow{PQ},\vec d_r,\vec d_s]=9$: \ \mbox{$d=\dfrac93=\res{3}$}


% =====================================================
\seccio{5}{Distància entre recta i pla i entre dos plans}

\ex{1} $\vec d\cdot\vec n=3-2-1=0$ i $P(2,\,0,\,1)\notin\pi$: paral·lels. \ \mbox{$d(r,\pi)=d(P,\pi)=\dfrac{|2-1+4|}{\sqrt6}=\res{\dfrac{5\sqrt6}{6}}$}

\ex{1} El segon és $3x-4y+12z+8=0$: \ \mbox{$d=\dfrac{|-5-8|}{\sqrt{9+16+144}}=\dfrac{13}{13}=\res{1}$}

\ex{2} $2x+y-2z+k=0$ amb \mbox{$\dfrac{|k|}{3}=3$}: \ $\res{2x+y-2z+9=0\ \text{i}\ 2x+y-2z-9=0}$

\ex{2} $\vec d\cdot\vec n=m+2-4=0$: \ $\res{m=2}$. \ $(1,\,1,\,1)\notin\pi$, \ \mbox{$d=\dfrac{|1+2-2|}{3}=\res{\dfrac13}$}

\ex{2} \pau{Andalusia, extraordinària 2025}\begin{tasks}(1)
  \task $\vec d=(1,\,-1,\,1)\times(1,\,2,\,-1)=(-1,\,2,\,3)$; \ $\vec d\cdot\vec n=-m-8=0$: \ $\res{m=-8}$ \ (i $r\not\subset\pi$)
  \task $P\left(\tfrac{10}{3},\,\tfrac13,\,0\right)\in r$: \ \mbox{$d=\dfrac{32}{\sqrt{69}}=\res{\dfrac{32\sqrt{69}}{69}}$}
\end{tasks}

\ex{2} \pau{Extremadura, ordinària 2025}\begin{tasks}(1)
  \task $\vec d\cdot\vec n=3-1-2=0$ i $(0,\,1,\,-1)\notin\pi$: \ \rest{paral·lels}
  \task \mbox{$d=\dfrac{|0+1+1-3|}{\sqrt3}=\res{\dfrac{\sqrt3}{3}}$}
  \task $\res{3x-y+2z+3=0}$
\end{tasks}

\ex{2} \pau{Galícia, extraordinària 2026}\begin{tasks}(1)
  \task $\overrightarrow{AQ}\times\vec d=(1,\,-2,\,1)\times(2,\,1,\,3)=(-7,\,-1,\,5)$: \ $\res{7x+y-5z-8=0}$
  \task $\vec d\cdot\vec n=2+1-3=0$; \ \mbox{$d=\dfrac{|1+1-0-1|}{\sqrt3}=\res{\dfrac{\sqrt3}{3}}$}
\end{tasks}

\ex{2} \pau{Múrcia, ordinària 2026}\begin{tasks}(1)
  \task $\vec d\cdot\vec n=4-1-3=0$ i $(1,\,-1,\,2)\notin\pi$: \ \rest{paral·lels}
  \task \mbox{$d=\dfrac{|2+1+6-5|}{\sqrt{14}}=\res{\dfrac{4}{\sqrt{14}}=\dfrac{2\sqrt{14}}{7}}$}
\end{tasks}

\ex{2} \pau{Castella i Lleó, extraordinària 2025}$r\parallel\pi$ ($-3+2+1=0$). $\pi':x+y-z+k=0$ amb \mbox{$\dfrac{|1+0-2+k|}{\sqrt3}=2\sqrt3$}, $|k-1|=6$: \ \rest{dos plans}, \ $\res{x+y-z+7=0}$ \ i \ $\res{x+y-z-5=0}$

\ex{2} \pau{Catalunya, extraordinària 2025}\begin{tasks}(1)
  \task \mbox{$\dfrac{|0-1+3-5|}{\sqrt6}=\dfrac{3}{\sqrt6}=\dfrac{\sqrt6}{2}$} \ \checkmark
  \task $\res{2x-y+z-2=0}$, \ distància \mbox{$\res{\dfrac{\sqrt6}{2}}$} (la de $P$)
  \task $2x-y+z+k=0$, $|k+5|=3$, $k\neq-2$: \ $\res{2x-y+z-8=0}$
\end{tasks}


% =====================================================
\seccio{6}{Recta que passa per un punt i talla dues rectes}

\ex{1} $\pi_1$ ($P$ i l'eix $OZ$): $x-2y=0$; \ $\pi_2$ ($P$ i $s$): $x-z-1=0$. \ $\res{t:\left\{\begin{aligned}x-2y&=0\\x-z-1&=0\end{aligned}\right.}$, \ director $(2,\,1,\,2)$. \ Talla $r$ en $\res{(0,\,0,\,-1)}$ i $s$ en $\res{\left(1,\,\tfrac12,\,0\right)}$

\ex{2} $\pi_1$ conté $r$ i $\vec v$: $y=0$; \ $\pi_2$ conté $s$ i $\vec v$: $x-z-1=0$. \ $\res{t:\left\{\begin{aligned}y&=0\\x-z-1&=0\end{aligned}\right.}$, \ és a dir, $(x,\,y,\,z)=(0,\,0,\,-1)+k\,(1,\,0,\,1)$, que talla $r$ en $(0,\,0,\,-1)$ i $s$ en $(1,\,0,\,0)$.

\ex{2} $r\cap\pi$: $3\lambda=3$, $X(1,\,2,\,0)$. La recta és $PX$: \ $\res{(x,\,y,\,z)=(1,\,1,\,1)+k\,(0,\,1,\,-1)}$ \ (està en $\pi$: $P\in\pi$ i $(0,\,1,\,-1)\perp(1,\,1,\,1)$).

\ex{2} $\pi_1$ ($O$ i $r$): $z=0$; \ $\pi_2$ ($O$ i $s$): $y=0$. La seua intersecció és l'eix $OX$, que talla $r$ en $(1,\,0,\,0)$ però \rest{és paral·lela a $s$} (i no la conté): no hi ha solució.

\ex{2} \pau{Navarra, extraordinària 2026}Pla per $A$ i $r$: normal $(1,\,0,\,3)\times(3,\,6,\,-3)\parallel(3,\,-2,\,-1)$, \ $3x-2y-z=0$. \ Pla per $A$ i $s$: normal $(-1,\,2,\,1)\times(0,\,5,\,-2)\parallel(9,\,2,\,5)$, \ $9x+2y+5z-16=0$. \ Director $(1,\,3,\,-3)$: \ \mbox{$\res{t:\dfrac{x-2}{1}=\dfrac{y-4}{3}=\dfrac{z+2}{-3}}$}. Talla $r$ en $(0,\,-2,\,4)$ i $s$ en $(1,\,1,\,1)$.

\ex{2} \pau{Castella i Lleó, extraordinària 2026}$s\cap\pi$: $(t,\,t,\,t)$ amb $2t=0$: $O(0,\,0,\,0)$. La recta és $AO$ ($A\in\pi$): \ $\res{r:\left\{\begin{aligned}x&=1+\lambda\\y&=-1-\lambda\\z&=0\end{aligned}\right.}$

\ex{2} \pau{Madrid, extraordinària 2024}\begin{tasks}(1)
  \task $\vec d_s=(1,\,-1,\,0)\times(1,\,0,\,1)=(-1,\,-1,\,1)$, $B(3,\,-2,\,0)$; \ $[\overrightarrow{AB},\vec d_r,\vec d_s]=4\neq0$: \rest{s'encreuen}; \ $|\vec d_r\times\vec d_s|=|(0,\,-4,\,-4)|=4\sqrt2$: \ \mbox{$d=\res{\dfrac{\sqrt2}{2}\approx0{,}71}$}
  \task Projecció de $P$ sobre $r$: pla $3x-y+z=18$, $M(5,\,-2,\,1)$. \ $\res{(x,\,y,\,z)=(5,\,-1,\,2)+k\,(0,\,1,\,1)}$
\end{tasks}


% =====================================================
\seccio{7}{Perpendicular comuna}

\ex{2} $X(\lambda,\,\lambda,\,0)$, $Y(1,\,\mu,\,\mu)$: \ $1-2\lambda+\mu=0$, \ $2\mu-\lambda=0$ \ $\Rightarrow$ \ $\lambda=\tfrac23$, $\mu=\tfrac13$: \ $\res{X\left(\tfrac23,\,\tfrac23,\,0\right),\ Y\left(1,\,\tfrac13,\,\tfrac13\right)}$; \ perpendicular comuna: \ $\res{(x,\,y,\,z)=\left(\tfrac23,\,\tfrac23,\,0\right)+t\,(1,\,-1,\,1)}$. \ $|\overrightarrow{XY}|=\left|\left(\tfrac13,\,-\tfrac13,\,\tfrac13\right)\right|=\tfrac{\sqrt3}{3}$ \ \checkmark

\ex{2} $\vec w\parallel(0,\,0,\,1)$; \ $\pi_r:x-y-1=0$, \ $\pi_s:x+y=0$: \ $\res{t:\left\{\begin{aligned}x-y-1&=0\\x+y&=0\end{aligned}\right.}$, \ és a dir, $x=\tfrac12$, $y=-\tfrac12$. \ Punts genèrics: $\lambda=-\tfrac12$, $\mu=\tfrac12$, \ $X\left(\tfrac12,\,-\tfrac12,\,0\right)$, $Y\left(\tfrac12,\,-\tfrac12,\,1\right)$; \ $d=|\overrightarrow{XY}|=\res{1}$

\ex{2} \pau{Múrcia, ordinària 2025}\begin{tasks}(1)
  \task $[\overrightarrow{PQ},\vec d_r,\vec d_s]=\begin{vmatrix}6 & -10 & 7\\0 & 1 & 1\\3 & 1 & 0\end{vmatrix}=-57\neq0$
  \task $\vec n=\vec d_r\times\vec d_s=(-1,\,3,\,-3)$; plans que contenen cada recta i $\vec n$: \ $\res{t:\left\{\begin{aligned}6x+y-z-11&=0\\3x-9y-10z+4&=0\end{aligned}\right.}$
  \task \mbox{$d=\dfrac{57}{\sqrt{19}}=\res{3\sqrt{19}}$}
\end{tasks}

\ex{2} \pau{Navarra, ordinària 2026}$X(\lambda,\,-3+\lambda,\,2-2\lambda)$, $Y(-4+3\mu,\,1+\mu,\,-1+2\mu)$; \ $\overrightarrow{XY}\perp\vec d_r,\vec d_s$ \ $\Rightarrow$ \ $X(1,\,-2,\,0)$, $Y(-1,\,2,\,1)$: \ \mbox{$\res{t:\dfrac{x+1}{2}=\dfrac{y-2}{-4}=\dfrac{z-1}{-1}}$}

\ex{2} \pau{Navarra, extraordinària 2025}$\vec d_r=(1,\,1,\,-2)\times(1,\,0,\,1)=(1,\,-3,\,-1)$, $P(3,\,7,\,0)\in r$. \ Punts genèrics: $X(3+\lambda,\,7-3\lambda,\,-\lambda)$, $Y(4-\mu,\,-1+\mu,\,-\mu)$; \ $\overrightarrow{XY}\perp\vec d_r,\vec d_s$ \ $\Rightarrow$ \ $X(5,\,1,\,-2)$, $Y(3,\,0,\,-1)$: \ \mbox{$\res{t:\dfrac{x-3}{2}=\dfrac{y}{1}=\dfrac{z+1}{-1}}$}


% =====================================================
\seccio{8}{Simètrics respecte d'un punt i d'una recta}

\ex{1} $A'=2B-A=\res{(-1,\,4,\,-5)}$. \qquad $M$ és el punt mig: \ \mbox{$M=\dfrac{A+A'}{2}=\res{(3,\,-1,\,2)}$}

\ex{2} Pla: $x+z=2$; \ tall: $\lambda+1+\lambda=2$, $\lambda=\tfrac12$, \ $M\left(\tfrac12,\,0,\,\tfrac32\right)$; \ $P'=2M-P=\res{(-1,\,-1,\,3)}$

\ex{2} \begin{tasks}(1)
  \task $M(1,\,0,\,0)$: \ $\res{P'(1,\,-2,\,-3)}$
  \task $\vec d=(1,\,-1,\,0)$, $A(1,\,1,\,0)$; \ $\overrightarrow{AP}\cdot\vec d=0$, de manera que $M=A=(1,\,1,\,0)$: \ $\res{P'(2,\,2,\,-1)}$
\end{tasks}

\ex{2} \pau{Castella i Lleó, ordinària 2024}Pla per $P$ perpendicular: $x+2y+2z+1=0$; \ $9\lambda+2=0$, \ $M(\tfrac{7}{9},\,-\tfrac{4}{9},\,-\tfrac{4}{9})$: \ $\res{P'(\tfrac{5}{9},\,-\tfrac{8}{9},\,\tfrac{1}{9})}$

\ex{2} \pau{País Basc, extraordinària 2025}$\vec d=(0,\,-1,\,-1)$; pla per $P$: $y+z+3=0$; \ $M(-2,\,-2,\,-1)$: \ $\res{P'(-8,\,-1,\,-2)}$


% =====================================================
\seccio{9}{Simètrics respecte d'un pla}

\ex{1} $(3+2t,\ 1-t,\ 2+t)$: \ $6+6t=0$, $t=-1$, \ $M(1,\,2,\,1)$; \ $P'=\res{(-1,\,3,\,0)}$

\ex{1} \begin{tasks}(1)
  \task $\res{(1,\,2,\,-3)}$
  \task $x-y=0$: \ $M\left(\tfrac32,\,\tfrac32,\,3\right)$, \ $\res{(2,\,1,\,3)}$
\end{tasks}

\ex{2} És perpendicular a $\overrightarrow{PP'}=(2,\,-2,\,2)$ i passa pel punt mig $(2,\,1,\,0)$: \ $\res{x-y+z-1=0}$

\ex{3} \begin{tasks}(1)
  \task Tall: $3+3\lambda=0$, \ $X(0,\,-1,\,1)$. \ Simètric de $A(1,\,1,\,1)$: $M=O$, $A'(-1,\,-1,\,-1)$. \ $r':\ \res{(x,\,y,\,z)=(0,\,-1,\,1)+\mu\,(1,\,0,\,2)}$
  \task És la recta que passa per $X$ i per la projecció $M=O$ de $A$: \ $\res{(x,\,y,\,z)=\mu\,(0,\,-1,\,1)}$
\end{tasks}

\ex{2} $\vec d\cdot\vec n=0$ i $(2,\,0,\,0)\notin\pi$: paral·lels. \ Simètric de $A(2,\,0,\,0)$: $M(2,\,2,\,2)$, $A'(2,\,4,\,4)$. \ $r':\ \res{(x,\,y,\,z)=(2,\,4,\,4)+\mu\,(0,\,1,\,-1)}$

\ex{2} \pau{Andalusia, ordinària 2024}\begin{tasks}(1)
  \task $\vec d=(1,\,1,\,1)$, $A(4,\,1,\,0)$; pla per $P$: $x+y+z=5$, que talla $r$ en $M=A$: \ $\res{P'(6,\,0,\,-1)}$
  \task $(1+t,\ -1-2t,\ -3+t)$: $6+6t=0$, $M(0,\,1,\,-4)$: \ $\res{Q'(-1,\,3,\,-5)}$
\end{tasks}

\ex{2} \pau{Andalusia, extraordinària 2026}\begin{tasks}(1)
  \task $r\cap\pi=(-1,\,0,\,4)$; director $\vec n\times\vec d_r\parallel(3,\,2,\,1)$: \ $\res{(x,\,y,\,z)=(-1,\,0,\,4)+\lambda\,(3,\,2,\,1)}$
  \task $\res{P'(2,\,-1,\,2)}$
\end{tasks}

\ex{2} \pau{Canàries, ordinària 2025}\begin{tasks}(1)
  \task $\vec d\parallel(3,\,2,\,1)$, $\vec d\cdot\vec n=2\neq0$: \ \rest{es tallen en un punt}
  \task $\res{A'\left(\tfrac53,\,-\tfrac43,\,\tfrac43\right)}$
\end{tasks}

\ex{2} \pau{Extremadura, extraordinària 2026}\begin{tasks}(1)
  \task $2-k=0$: \ $\res{k=2}$
  \task $(2+t,\ 3+2t,\ -4-3t)$: $14+14t=0$, $M(1,\,1,\,-1)$: \ $\res{A'(0,\,-1,\,2)}$
  \task \mbox{$\dfrac{14}{\sqrt{14}}=\res{\sqrt{14}}$}
\end{tasks}

\ex{2} \pau{País Basc, extraordinària 2026}$1+t=s$, $0=-1+s$, $-t=0$: \ $P(1,\,0,\,0)$. \ $(1+2t,\ t,\ -t)$: $6t-1=0$: \ $\res{P'\left(\tfrac53,\,\tfrac13,\,-\tfrac13\right)}$

\ex{2} \pau{Madrid, ordinària 2026}\begin{tasks}(1)
  \task $(2+2t,\ 2t,\ -t)$: $9t-9=0$, $M(4,\,2,\,-1)$: \ $\res{P'(6,\,4,\,-2)}$
  \task $\vec d\cdot\vec n=2+2-4=0$: paral·lels; \ \mbox{$d=\dfrac{|4-13|}{3}=\res{3}$}
\end{tasks}

\ex{2} \pau{Galícia, extraordinària 2026}\begin{tasks}(1)
  \task $P'=2Q-P=\res{(1,\,-2,\,5)}$
  \task $(3+t,\ 2t,\ -1+t)$: $6t-6=0$, $M(4,\,2,\,0)$: \ $\res{P''(5,\,4,\,1)}$
\end{tasks}


% =====================================================
\seccio{10}{Àrees i volums}

\ex{1} \begin{tasks}(1)
  \task $\overrightarrow{AB}\times\overrightarrow{AC}=(2,\,-2,\,1)\times(-1,\,-1,\,2)=(-3,\,-5,\,-4)$: \ \mbox{$S=\dfrac{\sqrt{50}}{2}=\res{\dfrac{5\sqrt2}{2}}$}
  \task \mbox{$\cos\hat A=\dfrac{-2+2+2}{3\sqrt6}=\dfrac{\sqrt6}{9}$} \ $\Rightarrow$ \ $\res{\hat A\approx74{,}21^\circ}$
  \task \mbox{$h=\dfrac{2S}{|\overrightarrow{AB}|}=\res{\dfrac{5\sqrt2}{3}}$}
\end{tasks}

\ex{2} $[\overrightarrow{AB},\overrightarrow{AC},\overrightarrow{AD}]=\begin{vmatrix}-1 & 1 & 0\\-1 & 0 & 1\\1 & 2 & 2\end{vmatrix}=5$: \ $\res{V=\tfrac56}$. \ Pla $ABC$: $x+y+z-1=0$; \ \mbox{$h=d(D,\pi)=\dfrac{5}{\sqrt3}=\res{\dfrac{5\sqrt3}{3}}$}

\ex{2} $V=\tfrac16\begin{vmatrix}1 & 0 & 0\\0 & 1 & 0\\1 & 1 & a\end{vmatrix}=\dfrac{|a|}{6}=2$ \ $\Rightarrow$ \ $\res{a=\pm12}$

\ex{2} $D=A+C-B=\res{(3,\,1,\,2)}$. \ $S=|\overrightarrow{AB}\times\overrightarrow{AD}|=|(1,\,1,\,-1)\times(2,\,1,\,1)|=|(2,\,-3,\,-1)|=\res{\sqrt{14}}$

\ex{2} $A(6,\,0,\,0)$, $B(0,\,4,\,0)$, $C(0,\,0,\,2)$: \ $V=\tfrac16\cdot6\cdot4\cdot2=\res{8}$. \ \mbox{$d(O,\pi)=\dfrac{12}{7}$}. \ \mbox{$S=\dfrac{3V}{h}=\dfrac{24}{12/7}=\res{14}$}

\ex{2} \pau{Cantàbria, extraordinària 2026}\begin{tasks}(1)
  \task $\overrightarrow{AB}=(-3,\,6,\,-9)=-3\,(1,\,-2,\,3)$ i $A\notin r$: \ \rest{paral·leles}
  \task Base $|\overrightarrow{AB}|=3\sqrt{14}$, altura \mbox{$d(A,r)=\dfrac{\sqrt{426}}{\sqrt{14}}$} (constant): \ \mbox{$S=\res{\dfrac{3\sqrt{426}}{2}\approx30{,}96}$}
  \task $\res{19x+8y-z-41=0}$
\end{tasks}

\ex{2} \pau{Cantàbria, ordinària 2024}\begin{tasks}(1)
  \task Amb els vectors que ixen de cada vèrtex: \ $\res{\hat A\approx43{,}05^\circ,\ \hat B\approx83{,}07^\circ,\ \hat C\approx53{,}88^\circ}$
  \task $\overrightarrow{AB}\times\overrightarrow{AC}=(-3,\,-2,\,6)\times(-8,\,-1,\,3)$, de mòdul $\sqrt{1690}$: \ \mbox{$S=\res{\dfrac{13\sqrt{10}}{2}}$}
\end{tasks}

\ex{2} \pau{Extremadura, ordinària 2024}\begin{tasks}(1)
  \task $\overrightarrow{AB}\times\overrightarrow{AC}=(-2,\,-2,\,2)$: \ $\res{\left(-\tfrac{1}{\sqrt3},\,-\tfrac{1}{\sqrt3},\,\tfrac{1}{\sqrt3}\right)}$ (o l'oposat)
  \task \mbox{$\cos\alpha=\dfrac{-2}{\sqrt2\,\sqrt8}=-\dfrac12$}: \ $\res{120^\circ}$
  \task $S=\tfrac12\sqrt{12}=\res{\sqrt3}$
\end{tasks}

\ex{2} \pau{Madrid, ordinària 2025}\begin{tasks}(1)
  \task \mbox{$\res{\dfrac{\sqrt2}{2}}$}
  \task És la projecció: \ $\res{Q\left(\tfrac12,\,\tfrac32,\,1\right)}$
  \task $\pi$ talla els eixos en $(2,\,0,\,0)$ i $(0,\,2,\,0)$ (és paral·lel a $OZ$): \ $S=\tfrac12|(2,\,-1,\,-1)\times(0,\,1,\,-1)|=\tfrac12|(2,\,2,\,2)|=\res{\sqrt3}$
\end{tasks}

\ex{2} \pau{Navarra, extraordinària 2024}$\vec d=(1,\,-2,\,2)$; la projecció de $P$ sobre $r$ és $(0,\,1,\,-2)$, el punt mig del segon costat, i el costat mesura $d(P,r)=6$. Els vèrtexs disten 3 d'aquest punt: $(0,\,1,\,-2)\pm(1,\,-2,\,2)$: \ $\res{(1,\,-1,\,0)\ \text{i}\ (-1,\,3,\,-4)}$


% =====================================================
\seccio{11}{Llocs geomètrics}

\ex{1} $(x-2)^2+(y-1)^2+(z+1)^2=x^2+(y-3)^2+(z-1)^2$ \ $\Rightarrow$ \ $-4x+4y+4z-4=0$. \ O bé: $\overrightarrow{AB}=(-2,\,2,\,2)$, $M(1,\,2,\,0)$. \ $\res{x-y-z+1=0}$

\ex{2} $|\vec n_1|=|\vec n_2|=\sqrt3$: \ $x+y-z-1=\pm(x-y+z-3)$: \ $\res{y-z+1=0}$ \ i \ $\res{x=2}$ \ (perpendiculars entre si).

\ex{1} $|2x-y+2z-1|=|2x-y+2z-7|$: el signe $+$ és impossible i el $-$ dona \ $\res{2x-y+2z=4}$, el pla paral·lel situat a mitja distància.

\ex{2} $R=|\overrightarrow{CP}|=|(2,\,2,\,1)|=3$: \ $\res{(x-1)^2+y^2+(z+2)^2=9}$. \ El pla tangent és perpendicular al radi $\overrightarrow{CP}$: \ $\res{2x+2y+z-9=0}$

\ex{2} \mbox{$R=d(C,\pi)=\dfrac{|2-2+2+7|}{3}=3$}: \ $\res{(x-2)^2+(y+1)^2+(z-1)^2=9}$. \ Punt de tangència = projecció de $C$: \ $\res{(1,\,-3,\,-1)}$

\ex{3} Ha d'estar en el pla mediador de $AB$ ($x=y$) i en el de $BC$ ($y=z$): \ és la recta \ $\res{x=y=z}$, perpendicular al pla $ABC$ pel seu circumcentre.

\ex{2} \pau{Andalusia, extraordinària 2024}$C(x,\,0,\,z)$ amb $|\overrightarrow{CA}|^2=|\overrightarrow{CB}|^2=|\overrightarrow{AB}|^2=20$: \ $(x-4)^2+z^2=x^2+4+z^2=20$ \ $\Rightarrow$ \ $x=\tfrac32$, $z^2=\tfrac{55}{4}$: \ $\res{C\left(\tfrac32,\,0,\,\pm\tfrac{\sqrt{55}}{2}\right)}$

\ex{2} \pau{Andalusia, ordinària 2024}\begin{tasks}(1)
  \task $\vec d_r=(1,\,0,\,2)$, $\vec d_s=(1,\,-1,\,0)$, $[\overrightarrow{OQ},\vec d_r,\vec d_s]\neq0$ amb $Q(-7,\,0,\,0)$: \ \rest{s'encreuen}
  \task $\vec n=\vec d_r\times\vec d_s=(2,\,2,\,-1)$; \ $2x+2y-z+k=0$ amb $|k|=|k-14|$: \ $\res{2x+2y-z+7=0}$
\end{tasks}

\ex{2} \pau{Castella i Lleó, ordinària 2026}\begin{tasks}(1)
  \task $\vec v_1\times\vec v_2=(2b,\,6,\,-2)\parallel(2,\,-3,\,1)$: $\res{b=-2}$; aleshores $\pi_2:2x-3y+z=-2$, i són distints si $\res{a\neq-2}$
  \task $\pi_2:3y-z=2$: \ $\res{x=\tfrac32,\ y=\lambda,\ z=-2+3\lambda}$
  \task Plans paral·lels $2x-3y+z=4$ i $2x-3y+z=-2$: el pla a mitja distància, \ $\res{2x-3y+z-1=0}$
\end{tasks}

\ex{2} \pau{Catalunya, ordinària 2025}\begin{tasks}(1)
  \task $\vec n'\parallel\overrightarrow{PQ}\times\vec n=(2,\,-2,\,4)\times(1,\,1,\,0)$: \ $\res{x-y-z=0}$
  \task És $\pi'\cap$ (pla mediador $x-y+2z=12$): \ $\res{x=2+\lambda,\ y=-2+\lambda,\ z=4}$
\end{tasks}

\ex{2} \pau{Catalunya, ordinària 2024}\begin{tasks}(1)
  \task $\res{x+y+1=0}$; \ $d(P,A)^2=d(P,B)^2$ es redueix a $8x+8y+8=0$.
  \task $d(A,\pi)=d(B,\pi)=\res{2\sqrt2}$: no és casualitat, $B$ és el simètric de $A$ respecte de $\pi$.
  \task $C\in\pi$, per tant $CA=CB$: \ \rest{isòsceles}. \ $S=\tfrac12|(-4,\,-4,\,0)\times(-8,\,4,\,0)|=\res{24}$
\end{tasks}

\ex{2} \pau{Múrcia, extraordinària 2024}\begin{tasks}(1)
  \task \mbox{$\cos\alpha=\dfrac13$}: \ $\res{\alpha\approx70{,}53^\circ}$
  \task $\vec d\cdot\vec n=2+1-3=0$ i $(1,\,-1,\,0)$ no és del pla: \ \rest{paral·lels}
  \task $(1+2t,\,-1+t,\,3t)$: \mbox{$\dfrac{|6t+3|}{\sqrt3}=\dfrac{3}{\sqrt3}$}, $t=0$ o $t=-1$: \ $\res{(1,\,-1,\,0)\ \text{i}\ (-1,\,-2,\,-3)}$
\end{tasks}

\ex{2} \pau{Navarra, extraordinària 2025}Pla mediador: \ $\res{x+2y+2z-10=0}$. \ $x+2y+2z+k=0$ amb \mbox{$\dfrac{|k+10|}{3}=6$}: \ $\res{x+2y+2z+8=0}$ \ i \ $\res{x+2y+2z-28=0}$

\ex{2} \pau{Madrid, extraordinària 2024}\begin{tasks}(1)
  \task Pla mediador: \ $\res{2x+4y-8z+5=0}$
  \task $R(2+t,\,t,\,t)$: $21+|PR|^2+|QR|^2=34+6t^2-2t=34$, \ $t=0$ (no val: $R(2,\,0,\,0)$) o $t=\tfrac13$: \ $\res{R\left(\tfrac73,\,\tfrac13,\,\tfrac13\right)}$
\end{tasks}

\ex{2} \pau{País Basc, ordinària 2026}\begin{tasks}(1)
  \task $\res{2x+y-z+2=0}$
  \task \mbox{$\dfrac{|2+1-2+2|}{\sqrt6}=\res{\dfrac{\sqrt6}{2}}$}
\end{tasks}


% =====================================================
\seccio{12}{Problemes i qüestions}

\ex{2} $(1+t,\ 2,\ 5-t)$ amb $z=0$: \ $t=5$, \ $\res{X(6,\,2,\,0)}$. \ $P'(1,\,2,\,-5)$, \ $\overrightarrow{P'X}=(5,\,0,\,5)$: \ raig reflectit \ $\res{(x,\,y,\,z)=(6,\,2,\,0)+\mu\,(1,\,0,\,1)},\ \mu\ge0$

\ex{3} Si $B'(4,\,5,\,-1)$ és el simètric de $B$, $d(X,B)=d(X,B')$ i el camí més curt és el segment $AB'$. \ $A+t\,(3,\,4,\,-3)$ amb $z=0$: $t=\tfrac23$: \ $\res{X\left(3,\,\tfrac{11}{3},\,0\right)}$; \ suma mínima \ $|\overrightarrow{AB'}|=\res{\sqrt{34}}$

\ex{2} \begin{tasks}(1)
  \task \rest{Falsa}: per exemple, $x=0$ i $y=0$ són perpendiculars a $z=0$ i es tallen.
  \task \rest{Certa}: el seu director és proporcional a $\vec n$, que és perpendicular a qualsevol vector del pla (també si la recta del pla no passa pel punt de tall: s'encreuen perpendicularment).
  \task \rest{Falsa}: si s'encreuen no tenen cap punt comú i $[\overrightarrow{PQ},\vec d_r,\vec d_s]\neq0$, per tant $d>0$.
  \task \rest{Certa}: $|Ax_0+By_0+Cz_0+D|=0$ vol dir que $P$ compleix l'equació del pla.
\end{tasks}

\ex{2} \pau{Astúries, ordinària 2026}\begin{tasks}(1)
  \task Punt mig: \ $\res{G(1,\,0,\,2)}$
  \task $\overrightarrow{GH}=(0,\,1,\,-1)$: \ $\res{x=1,\ y=\lambda,\ z=2-\lambda}$
  \task $\overrightarrow{PQ}\cdot\overrightarrow{GH}=(-2,\,-2,\,-2)\cdot(0,\,1,\,-1)=0$: \ $\res{90^\circ}$
\end{tasks}

\ex{2} \pau{Múrcia, extraordinària 2025}\begin{tasks}(1)
  \task \mbox{$\res{\dfrac{x-1}{1}=\dfrac{y-2}{1}=\dfrac{z-1}{3}}$}
  \task $6+11t=0$: \ $\res{\left(\tfrac{5}{11},\,\tfrac{16}{11},\,-\tfrac{7}{11}\right)}$
  \task \mbox{$d=\dfrac{6}{\sqrt{11}}=\res{\dfrac{6\sqrt{11}}{11}\approx1{,}81}$}
\end{tasks}

\ex{2} \pau{Astúries, ordinària 2025}\begin{tasks}(1)
  \task $[\overrightarrow{PQ},\vec d_1,\vec d_2]=27\neq0$: \ \rest{s'encreuen}
  \task Perpendicular comuna: \ $\res{E(1,\,1,\,1),\ F(3,\,3,\,2)}$, \ longitud $|\overrightarrow{EF}|=\res{3}$
\end{tasks}

\ex{2} \pau{Madrid, extraordinària 2026}\begin{tasks}(1)
  \task Costat \mbox{$=d(P,\pi)=\dfrac{|2-4+1-1|}{3}=\res{\dfrac23}$}; \ volum \mbox{$\res{\dfrac{8}{27}}$}
  \task El vèrtex de la cara del pla que està sota $P$ (la projecció): \ $\res{\left(\tfrac{13}{9},\,\tfrac{14}{9},\,\tfrac{11}{9}\right)}$
\end{tasks}

\ex{3} \pau{Aragó, extraordinària 2025}\begin{tasks}(1)
  \task $D(3\lambda,\,1+\lambda,\,3)$; \ \mbox{$V=\tfrac16\big|[\overrightarrow{AB},\overrightarrow{AC},\overrightarrow{AD}]\big|=\dfrac{67-23\lambda}{24}$}, màxim en $\lambda=0$: \ \mbox{$\res{V=\dfrac{67}{24}\approx2{,}79\ \text{m}^3}$}
  \task \mbox{$V=\dfrac{11}{6}$}, \ \mbox{$S_{ABC}=\tfrac12\left|\left(-\tfrac32,\,-\tfrac54,\,6\right)\right|=\dfrac{7\sqrt{13}}{8}$}: \ \mbox{$h=\dfrac{3V}{S}=\res{\dfrac{44\sqrt{13}}{91}\approx1{,}74\ \text{m}}$}
\end{tasks}

\ex{2} \pau{Castella-la Manxa, ordinària 2026}\begin{tasks}(1)
  \task $M(0,\,1,\,5)$: \ $\res{P'(-1,\,0,\,4)}$
  \task $|\overrightarrow{PP'}|=2\,d(P,\pi)=\res{2\sqrt3}$
\end{tasks}

\ex{2} \pau{Castella-la Manxa, extraordinària 2025}\begin{tasks}(1)
  \task $[\overrightarrow{AB},\overrightarrow{AC},\overrightarrow{AD}]=6-m=0$: \ $\res{m=6}$
  \task $\res{5x-3y-z+1=0}$
  \task \mbox{$d(P,\pi)=\dfrac{3}{\sqrt{35}}\approx0{,}51<1$}: \ \rest{no està ben situat}
\end{tasks}

\ex{2} \pau{Illes Balears, ordinària 2024}\begin{tasks}(1)
  \task $D=B+C-A=\res{(1,\,1,\,4)}$
  \task $\vec n=\overrightarrow{AB}\times\overrightarrow{AC}=(-7,\,-5,\,3)$, $|\vec n|=\sqrt{83}$: \ $\res{E=\pm\tfrac{1}{\sqrt{83}}(-7,\,-5,\,3)}$
  \task $\res{7x+5y-3z\pm\sqrt{83}=0}$
\end{tasks}


\end{document}
